\documentclass[10pt,a4paper]{amsart}
\usepackage{amsmath,amssymb,amsthm,graphicx}
\usepackage{verbatim}
\newenvironment{reference}{\comment}{\endcomment}
%\newenvironment{reference}{}{}
\newenvironment{slogan}{\comment}{\endcomment}
\newenvironment{history}{\comment}{\endcomment}
% For commutative diagrams we use Xy-pic
\usepackage[all]{xy}

% We use 2cell for 2-commutative diagrams.
\xyoption{2cell}
% Theorem environments.
%
\theoremstyle{plain}
\newtheorem{theorem}[subsection]{Theorem}
\newtheorem{proposition}[subsection]{Proposition}
\newtheorem{lemma}[subsection]{Lemma}

\theoremstyle{definition}
\newtheorem{definition}[subsection]{Definition}
\newtheorem{example}[subsection]{Example}
\newtheorem{exercise}[subsection]{Exercise}
\newtheorem{situation}[subsection]{Situation}

\theoremstyle{remark}
\newtheorem{remark}[subsection]{Remark}
\newtheorem{remarks}[subsection]{Remarks}

\numberwithin{equation}{subsection}

% Macros
%
\def\lim{\mathop{\mathrm{lim}}\nolimits}
\def\colim{\mathop{\mathrm{colim}}\nolimits}
\def\Spec{\mathop{\mathrm{Spec}}}
\def\Hom{\mathop{\mathrm{Hom}}\nolimits}
\def\Ext{\mathop{\mathrm{Ext}}\nolimits}
\def\SheafHom{\mathop{\mathcal{H}\!\mathit{om}}\nolimits}
\def\SheafExt{\mathop{\mathcal{E}\!\mathit{xt}}\nolimits}
\def\Sch{\mathit{Sch}}
\def\Mor{\mathop{\mathrm{Mor}}\nolimits}
\def\Ob{\mathop{\mathrm{Ob}}\nolimits}
\def\Sh{\mathop{\mathit{Sh}}\nolimits}
\def\NL{\mathop{N\!L}\nolimits}
\def\CH{\mathop{\mathrm{CH}}\nolimits}
\def\proetale{{pro\text{-}\acute{e}tale}}
\def\etale{{\acute{e}tale}}
\def\QCoh{\mathit{QCoh}}
\def\Ker{\mathop{\mathrm{Ker}}}
\def\Im{\mathop{\mathrm{Im}}}
\def\Coker{\mathop{\mathrm{Coker}}}
\def\Coim{\mathop{\mathrm{Coim}}}

% Boxtimes
%
\DeclareMathSymbol{\boxtimes}{\mathbin}{AMSa}{"02}

%
% Macros for moduli stacks/spaces
%
\def\QCohstack{\mathcal{QC}\!\mathit{oh}}
\def\Cohstack{\mathcal{C}\!\mathit{oh}}
\def\Spacesstack{\mathcal{S}\!\mathit{paces}}
\def\Quotfunctor{\mathrm{Quot}}
\def\Hilbfunctor{\mathrm{Hilb}}
\def\Curvesstack{\mathcal{C}\!\mathit{urves}}
\def\Polarizedstack{\mathcal{P}\!\mathit{olarized}}
\def\Complexesstack{\mathcal{C}\!\mathit{omplexes}}
% \Pic is the operator that assigns to X its picard group, usage \Pic(X)
% \Picardstack_{X/B} denotes the Picard stack of X over B
% \Picardfunctor_{X/B} denotes the Picard functor of X over B
\def\Pic{\mathop{\mathrm{Pic}}\nolimits}
\def\Picardstack{\mathcal{P}\!\mathit{ic}}
\def\Picardfunctor{\mathrm{Pic}}
\def\Deformationcategory{\mathcal{D}\!\mathit{ef}}
\usepackage{cleveref}
\crefname{originallemma}{original lemma}{original lemmas}
\Crefname{originallemma}{Original lemma}{Original lemmas}
\makeatletter
\newlabel{lemma-check-quasi-compact-covering}{{lemma-check-quasi-compact-covering}{source}{}{sourcepointer-1}{}}
\newlabel{lemma-check-quasi-compact-covering@cref}{{[originallemma][0][]lemma-check-quasi-compact-covering}{[1][0][]source}{}{}{}}
\makeatother
\crefname{originallemma}{original lemma}{original lemmas}
\Crefname{originallemma}{Original lemma}{Original lemmas}
\makeatletter
\newlabel{lemma-separated-over-separated}{{lemma-separated-over-separated}{source}{}{sourcepointer-2}{}}
\newlabel{lemma-separated-over-separated@cref}{{[originallemma][0][]lemma-separated-over-separated}{[1][0][]source}{}{}{}}
\makeatother
\crefname{originalreference}{original reference}{original references}
\Crefname{originalreference}{Original reference}{Original references}
\makeatletter
\newlabel{morphisms-lemma-finite-type-noetherian}{{morphisms-lemma-finite-type-noetherian}{source}{}{sourcepointer-3}{}}
\newlabel{morphisms-lemma-finite-type-noetherian@cref}{{[originalreference][0][]morphisms-lemma-finite-type-noetherian}{[1][0][]source}{}{}{}}
\makeatother
\crefname{originalreference}{original reference}{original references}
\Crefname{originalreference}{Original reference}{Original references}
\makeatletter
\newlabel{morphisms-lemma-immersion-locally-finite-type}{{morphisms-lemma-immersion-locally-finite-type}{source}{}{sourcepointer-4}{}}
\newlabel{morphisms-lemma-immersion-locally-finite-type@cref}{{[originalreference][0][]morphisms-lemma-immersion-locally-finite-type}{[1][0][]source}{}{}{}}
\makeatother
\crefname{originalreference}{original reference}{original references}
\Crefname{originalreference}{Original reference}{Original references}
\makeatletter
\newlabel{properties-lemma-immersion-into-noetherian}{{properties-lemma-immersion-into-noetherian}{source}{}{sourcepointer-5}{}}
\newlabel{properties-lemma-immersion-into-noetherian@cref}{{[originalreference][0][]properties-lemma-immersion-into-noetherian}{[1][0][]source}{}{}{}}
\makeatother
\crefname{originalreference}{original reference}{original references}
\Crefname{originalreference}{Original reference}{Original references}
\makeatletter
\newlabel{stacks-properties-definition-immersion}{{stacks-properties-definition-immersion}{source}{}{sourcepointer-6}{}}
\newlabel{stacks-properties-definition-immersion@cref}{{[originalreference][0][]stacks-properties-definition-immersion}{[1][0][]source}{}{}{}}
\makeatother
\title{Constructible locally closed subsets of the space of a Noetherian stack are open}
\author{The Stacks Project Authors}
\date{stacks-project:stacks-morphisms.tex:local-snapshot-20261012; GNU FDL 1.2 or later}
\begin{document}
\maketitle
\noindent\textit{Editorial note.} All delivered mathematical text is the author's own, copied from the pinned Stacks Project chapter source sources/stacks-algebraic-stacks-morphisms/pinned-original-source.tex; no statement or proof step was rewritten, reordered or removed.\par\medskip
\noindent\textit{Editorial note.} Preamble provenance: the Stacks Project keeps one single shared preamble for every chapter; this card takes it directly from the staged sources/stacks-algebraic-stacks-morphisms/preamble.tex as three verbatim spans (lines 8-12, 14-18 and 154-217), so no derived file was written and nothing was added or redefined. The wrapper class and the base packages are supplied by the assembler.\par\medskip
\noindent\textit{Editorial note.} Bibliography provenance: neither the counted statement nor its proof contains a citation (this lane rescans the whole span before writing the card), so the card carries no bibliography and no entry was reconstructed.\par\medskip
\noindent\textit{Editorial note.} Selection provenance: this card was opened under the file-based control contract (CONTROL.md), with the mechanical pre-card gate run on its statement and proof interval before the card was written.\par\medskip
\noindent\textit{Allowed standard prerequisites (editorial).}\begin{itemize}
\item The description of the topology of the space of an algebraic stack through a smooth presentation. Used at the start of the counted proof (source lines 2233-2240). Allows the statement to be tested on the presenting scheme.
\item The corresponding statement for schemes that a constructible subset of a Noetherian scheme is open. Invoked inside the counted proof (source lines 2241-2247) as a cross-chapter pointer. Gives openness for the presenting scheme.
\end{itemize}\medskip
\ifdefined\hypertarget\hypertarget{sourcepointer-1}{}\fi
\noindent\textit{Source reference.} Original lemma lemma-check-quasi-compact-covering: sources/stacks-algebraic-stacks-morphisms/pinned-original-source.tex, source lines 2154--2154. Navigation pointer only: the counted text refers to the Stacks item lemma-check-quasi-compact-covering, which is an existing labeled result of this same chapter and is not part of this card. The pointer carries the target's own label as its number and its definition line as its span; its content is not reproduced here.\par\medskip
\ifdefined\hypertarget\hypertarget{sourcepointer-2}{}\fi
\noindent\textit{Source reference.} Original lemma lemma-separated-over-separated: sources/stacks-algebraic-stacks-morphisms/pinned-original-source.tex, source lines 888--888. Navigation pointer only: the counted text refers to the Stacks item lemma-separated-over-separated, which is an existing labeled result of this same chapter and is not part of this card. The pointer carries the target's own label as its number and its definition line as its span; its content is not reproduced here.\par\medskip
\ifdefined\hypertarget\hypertarget{sourcepointer-3}{}\fi
\noindent\textit{Source reference.} Original reference morphisms-lemma-finite-type-noetherian: sources/stacks-algebraic-stacks-morphisms/pinned-original-source.tex, source lines 2240--2240. Navigation pointer only: the counted text refers to the Stacks item morphisms-lemma-finite-type-noetherian, which belongs to another chapter of the project and is not part of this card. The pointer carries the target's own label as its number because this chapter's staged snapshot does not contain the other chapter's tag table; its content is not reproduced here.\par\medskip
\ifdefined\hypertarget\hypertarget{sourcepointer-4}{}\fi
\noindent\textit{Source reference.} Original reference morphisms-lemma-immersion-locally-finite-type: sources/stacks-algebraic-stacks-morphisms/pinned-original-source.tex, source lines 2239--2239. Navigation pointer only: the counted text refers to the Stacks item morphisms-lemma-immersion-locally-finite-type, which belongs to another chapter of the project and is not part of this card. The pointer carries the target's own label as its number because this chapter's staged snapshot does not contain the other chapter's tag table; its content is not reproduced here.\par\medskip
\ifdefined\hypertarget\hypertarget{sourcepointer-5}{}\fi
\noindent\textit{Source reference.} Original reference properties-lemma-immersion-into-noetherian: sources/stacks-algebraic-stacks-morphisms/pinned-original-source.tex, source lines 2241--2241. Navigation pointer only: the counted text refers to the Stacks item properties-lemma-immersion-into-noetherian, which belongs to another chapter of the project and is not part of this card. The pointer carries the target's own label as its number because this chapter's staged snapshot does not contain the other chapter's tag table; its content is not reproduced here.\par\medskip
\ifdefined\hypertarget\hypertarget{sourcepointer-6}{}\fi
\noindent\textit{Source reference.} Original reference stacks-properties-definition-immersion: sources/stacks-algebraic-stacks-morphisms/pinned-original-source.tex, source lines 2236--2236. Navigation pointer only: the counted text refers to the Stacks item stacks-properties-definition-immersion, which belongs to another chapter of the project and is not part of this card. The pointer carries the target's own label as its number because this chapter's staged snapshot does not contain the other chapter's tag table; its content is not reproduced here.\par\medskip
\begin{lemma}
\label{lemma-locally-closed-in-noetherian}
Let $j : \mathcal{X} \to \mathcal{Y}$ be an immersion of algebraic stacks.
\begin{enumerate}
\item If $\mathcal{Y}$ is locally Noetherian, then
$\mathcal{X}$ is locally Noetherian and $j$ is quasi-compact.
\item If $\mathcal{Y}$ is Noetherian, then $\mathcal{X}$ is Noetherian.
\end{enumerate}
\end{lemma}

\begin{proof}
Choose a scheme $V$ and a surjective smooth morphism $V \to \mathcal{Y}$.
Then $U = \mathcal{X} \times_\mathcal{Y} V$ is a scheme and
$V \to U$ is an immersion, see
Properties of Stacks, Definition \ref{stacks-properties-definition-immersion}.
Recall that $\mathcal{Y}$ is locally Noetherian if and only if $V$
is locally Noetherian. In this case $U$ is locally Noetherian too
(Morphisms, Lemmas \ref{morphisms-lemma-immersion-locally-finite-type} and
\ref{morphisms-lemma-finite-type-noetherian}) and $U \to V$ is quasi-compact
(Properties, Lemma \ref{properties-lemma-immersion-into-noetherian}).
This shows that $j$ is quasi-compact
(Lemma \ref{lemma-check-quasi-compact-covering})
and that $\mathcal{X}$ is locally Noetherian.
Finally, if $\mathcal{Y}$ is Noetherian, then we see from the above
that $\mathcal{X}$ is quasi-compact and locally Noetherian.
To finish the proof observe that $j$ is separated and hence
$\mathcal{X}$ is quasi-separated because $\mathcal{Y}$ is so by
Lemma \ref{lemma-separated-over-separated}.
\end{proof}

\end{document}
